%% UML — Diagrama de secuencia (plantilla)
%% Compilar: tectonic -X compile uml-secuencia.tex
\documentclass[11pt,a4paper]{article}
\usepackage{../../docs/latex/legasoft}
\usepackage{../../docs/latex/legasoft-diagramas}

\lgtitulo{UML: diagrama de secuencia}
\lgtituloCorto{UML — Secuencia}
\lgcodigo{LGS-MOD-UML-3}
\lgversion{0.1}
\lgestado{Plantilla}
\lgfecha{\campo{fecha}}
\lgautor{\lgArquitecto\ (\lgArquitectoRol)}

\begin{document}

\begin{center}
{\sffamily\bfseries\LARGE\color{lgPrimario} UML: diagrama de secuencia\par}
\vspace{2pt}
{\sffamily\normalsize\color{lgGris} \campo{CU-01 Registrar información} · Legasoft\par}
\end{center}
\vspace{6pt}
{\color{lgBorde}\hrule height 0.8pt}
\vspace{10pt}

\begin{porcompletar}[Cómo usar esta plantilla]
Un diagrama de secuencia se elabora por escenario, normalmente el flujo principal
de un caso de uso. Indique en el título el caso de uso representado y mantenga la
correspondencia con los contratos documentados en \texttt{LGS-ARQ-001}.
\end{porcompletar}

% =============================================================================
\section*{Diagrama}

\begin{figure}[H]
\centering
\ajustaancho{%
\begin{tikzpicture}[uml, x=1cm, y=1cm]

  % --- Participantes --------------------------------------------------------
  \node[umlobjeto] (a) at (0,0)    {\campo{: Usuario}};
  \node[umlobjeto] (b) at (3.6,0)  {\campo{: AppWeb}};
  \node[umlobjeto] (c) at (7.2,0)  {\campo{: API}};
  \node[umlobjeto] (d) at (10.8,0) {\campo{: BaseDatos}};

  % --- Líneas de vida ------------------------------------------------------
  \def\vidafin{-8.2}
  \foreach \n in {a,b,c,d}{
    \draw[umlvida] (\n.south) -- ($(\n.south)+(0,\vidafin)$);}

  % --- Barras de activación ------------------------------------------------
  \fill[umlactivacion] ($(b.south)+(-0.10,-0.6)$)  rectangle ($(b.south)+(0.10,-7.4)$);
  \fill[umlactivacion] ($(c.south)+(-0.10,-1.3)$)  rectangle ($(c.south)+(0.10,-6.6)$);
  \fill[umlactivacion] ($(d.south)+(-0.10,-2.6)$)  rectangle ($(d.south)+(0.10,-3.9)$);
  \fill[umlactivacion] ($(d.south)+(-0.10,-4.6)$)  rectangle ($(d.south)+(0.10,-5.6)$);

  % --- Mensajes ------------------------------------------------------------
  \draw[umlmensaje] ($(a.south)+(0,-0.6)$)  -- ($(b.south)+(-0.10,-0.6)$)
    node[midway, above] {1: \campo{enviar formulario}};

  \draw[umlmensaje] ($(b.south)+(0.10,-1.3)$) -- ($(c.south)+(-0.10,-1.3)$)
    node[midway, above] {2: \campo{POST /registro}};

  \draw[umlmensaje] ($(c.south)+(0.10,-2.6)$) -- ($(d.south)+(-0.10,-2.6)$)
    node[midway, above] {3: \campo{validar unicidad}};

  \draw[umlretorno] ($(d.south)+(-0.10,-3.9)$) -- ($(c.south)+(0.10,-3.9)$)
    node[midway, above] {4: \campo{resultado}};

  \draw[umlmensaje] ($(c.south)+(0.10,-4.6)$) -- ($(d.south)+(-0.10,-4.6)$)
    node[midway, above] {5: \campo{insertar registro}};

  \draw[umlretorno] ($(d.south)+(-0.10,-5.6)$) -- ($(c.south)+(0.10,-5.6)$)
    node[midway, above] {6: \campo{id generado}};

  \draw[umlretorno] ($(c.south)+(-0.10,-6.6)$) -- ($(b.south)+(0.10,-6.6)$)
    node[midway, above] {7: \campo{201 Created}};

  \draw[umlretorno] ($(b.south)+(-0.10,-7.4)$) -- ($(a.south)+(0,-7.4)$)
    node[midway, above] {8: \campo{confirmación}};

\end{tikzpicture}}
\caption{Flujo principal del escenario: mensajes entre los participantes en el
tiempo. El eje vertical representa el avance temporal.}
\end{figure}

% =============================================================================
\Needspace*{20\baselineskip}
\section*{Escenario representado}

\begin{center}
\begin{tabularx}{\linewidth}{@{}L{3.6cm} Y@{}}
\toprule
\sffamily\bfseries\small\color{lgPrimario}Caso de uso     & \campo{CU-01 Registrar información} \\
\sffamily\bfseries\small\color{lgPrimario}Requisito       & \campo{RF-001} \\
\sffamily\bfseries\small\color{lgPrimario}Precondición    & \campo{el usuario está autenticado} \\
\sffamily\bfseries\small\color{lgPrimario}Postcondición   & \campo{el registro queda persistido y es consultable} \\
\sffamily\bfseries\small\color{lgPrimario}Flujos alternos & \campo{datos inválidos, registro duplicado, error de conexión} \\
\bottomrule
\end{tabularx}
\captionof{table}{Contexto del escenario representado.}
\end{center}

\begin{nota}[Notación]
La flecha continua con punta rellena representa un mensaje síncrono; la flecha
discontinua, el retorno. Las barras verticales sobre la línea de vida indican el
periodo en que el participante está activo. Los flujos alternos se modelan en
diagramas aparte o mediante fragmentos combinados (\texttt{alt}, \texttt{opt}).
\end{nota}

\end{document}
